\subsection{Perturbations of continuous injective or surjective linear maps}

In this section, we adopt the following conventions: if E is a finite-dimensional vector space, \(\dim(\mathrm{E})\) denotes its dimension; if E is an infinite-dimensional vector space, we set \(\dim(\mathrm{E}) = +\infty \in \overline{\mathbf{R}}\) . If u is a linear map whose kernel or cokernel is finite-dimensional, we set \(\operatorname{ind}(u) = \dim \operatorname{Coker}(u) - \dim \operatorname{Ker}(u)\) , the computation being carried out in \(\overline{R}\) .

Let E and F be normed spaces. We denote by \(\mathcal{M}(\mathrm{E};\mathrm{F})\) the set of strict injective morphisms from E into F, and by \(\mathcal{Q}\mathcal{M}(\mathrm{E};\mathrm{F})\) the set of strict morphisms from E into F whose kernel is finite-dimensional.

PROPOSITION 10. --- Let E and F be normed spaces. The set \(\mathcal{M}(\mathrm{E};\mathrm{F})\) is open in \(\mathcal{L}(\mathrm{E};\mathrm{F})\) . It is the interior of the set of maps in \(\mathcal{L}(\mathrm{E};\mathrm{F})\) that are injective.

Let A be the set of continuous injective linear maps from E into F. For a map \(u \in A\) to be a strict morphism, it is necessary and sufficient that its conorm \(((u))\) be strictly positive (III, p. 62, remark). Now \(((u))\) is the distance from \(u\) to the complement of A in \(\mathcal{L}(\mathrm{E};\mathrm{F})\) (III, p. 65, prop. 9). The proposition follows.

PROPOSITION 11. --- Let E and F be Banach spaces. The set \(\mathcal{Q}\mathcal{M}(\mathrm{E};\mathrm{F})\) is open in \(\mathcal{L}(\mathrm{E};\mathrm{F})\) . It is the interior of the set of maps in \(\mathcal{L}(\mathrm{E};\mathrm{F})\) whose kernel is finite-dimensional.

Let A be the set of continuous linear maps from E into F whose kernel is finite-dimensional.

Let u be an element of \(\mathcal{Q}\mathcal{M}(\mathrm{E};\mathrm{F})\) . We then have \(((u)) > 0\) (III, p. 62, remark). Every element \(v \in \mathcal{L}(\mathrm{E};\mathrm{F})\) whose distance from u is \(<((u))\) then belongs to A (III, p. 65, prop. 9), so that u is an interior point of A.

Conversely, let \(u\) be an element of A that is not a strict morphism. We have \(((u)) = 0\) (III, p. 62, remark). Let \(v\) be an element of \(\mathcal{L}(\mathrm{E};\mathrm{F})\) which differs from \(u\) by a linear map of finite rank; by Corollary 1 of III, p. 40, the morphism \(v\) cannot be strict with closed image; since a strict morphism from E into F has a closed image in F (EVT, IV, p. 28, th. 1), this means that v is not a strict morphism. We therefore have \(((v)) = 0\) . Let \(\varepsilon\) be a real number \textgreater{} 0. By the preceding observations, Lemma 3 of III, p. 66 allows us to construct by induction a sequence \((u_{m})_{m \in \mathbf{N}}\) of elements of \(\mathcal{L}(\mathrm{E}; \mathrm{F})\) satisfying the following conditions:

\begin{enumerate}
\def\labelenumi{(\roman{enumi})}
\item
  One has \(u_{0}=u\) ;
\item
  For every \(m \geqslant 0\) , \(u_{m+1} - u_m\) is a continuous linear map of rank 1 and norm \(\leqslant 2^{-m-1}\varepsilon\) ;
\item
  For every \(m \geqslant 0\) , the kernel of \(u_m\) has dimension \(n + m\) and is contained in that of \(u_{m+1}\) .
\end{enumerate}

The sequence \((u_{m})\) is a Cauchy sequence in the Banach space \(\mathcal{L}(\mathrm{E};\mathrm{F})\) . Let \(u'\) be its limit. The kernel of \(u'\) contains that of \(u_{m}\) for all \(m \geqslant 0\) ; it is of infinite dimension. Since

\[
\left\| u ^ {\prime} - u \right\| \leqslant \sum_ {m = 0} ^ {\infty} \left\| u _ {m + 1} - u _ {m} \right\| \leqslant \varepsilon \sum_ {m = 0} ^ {\infty} 2 ^ {- m - 1} = \varepsilon ,
\]

the distance from u to the complement of A is less than \(\varepsilon\) . Since this holds for every \(\varepsilon > 0\) , one concludes that u is not an interior point of A.

PROPOSITION 12. --- Let E and F be Banach spaces and u an element of \(\mathcal{Q}\mathcal{M}(\mathrm{E};\mathrm{F})\) . Every \(v \in \mathcal{L}(\mathrm{E};\mathrm{F})\) such that \(\|v - u\| < ((u))\) belongs to \(\mathcal{Q}\mathcal{M}(\mathrm{E};\mathrm{F})\) and satisfies the relations

\[
\begin{array}{c} \dim \operatorname{Ker} (v) \leqslant \dim \operatorname{Ker} (u), \\ \dim \operatorname{Coker} (v) \leqslant \dim \operatorname{Coker} (u), \\ \operatorname{ind} (v) = \operatorname{ind} (u). \end{array}
\]

Since \(u\) is strict, we have \(((u)) > 0\) . Let B denote the open ball with center \(v\) and radius \(((u))\) in \(\mathcal{L}(\mathrm{E};\mathrm{F})\) . For every \(v \in \mathrm{B}\) , we have \(\dim \operatorname{Ker}(v) \leqslant \dim \operatorname{Ker}(u)\) (III, p. 65, prop. 9) and \(v \in \mathcal{Q}\mathcal{M}(\mathrm{E};\mathrm{F})\) (prop. 11).

For \(r \in \mathbf{Z} \cup \{+\infty\}\) , denote by \(B_r\) the set of elements \(v \in B\) such that \(\text{ind}(v) = r\) . If \(r \in \mathbf{Z}\) , the set \(B_r\) is the set of Fredholm maps from E to F of index \(r\) which belong to B (III, p. 52, prop. 11); it is open in B by th. 1 of III, p. 58.

Let us prove that the sets \(B_{r}\) are closed in B. Let \(v \in B\) be a point adherent to \(B_{r}\) . Since the sets \(B_{s}\) , for \(s \in Z\) , are open in B and pairwise disjoint, we have \(v \in B_{r}\) or \(v \in B_{+\infty}\) .

Suppose \(v \in B_{+\infty}\) . Let n denote the dimension of \(\operatorname{Ker}(u)\) . Choose a vector subspace T of F of dimension \(r + 2n + 1\) whose intersection with \(\operatorname{Im}(v)\) is reduced to 0 and a topological supplement S of \(\operatorname{Ker}(v)\) in E (III, p. 55, prop. 1). Consider the continuous linear map \(f: (s, t) \mapsto v(s) + t\) from S × T to F. It is injective. The linear map v is a strict morphism since v belongs to B ⊂ \(\mathcal{Q}\mathcal{M}(\mathrm{E};\mathrm{F})\) . The image of v is therefore closed (EVT, IV, p. 28, th. 1). By prop. 1 of III, p. 39, the restriction of v to S is a strict morphism with closed image from S to F, and f is a strict morphism with closed image from S × T to F.

By prop. 10, there exists a neighborhood U of \(v\) in B such that, for every \(w \in \mathrm{U}\) , the linear map \((s,t) \mapsto w(s) + t\) from \(\mathrm{S} \times \mathrm{T}\) into F is injective. Let \(w\) be an element of U. We then have \(\operatorname{Ker}(w) \cap \mathrm{S} = \{0\}\) , hence \(w\) induces by restriction and passage to the quotient an injective linear map from \(\operatorname{Ker}(w)\) into E/S, which implies that \(\operatorname{Ker}(w)\) has dimension at most \(n\) . We also have \(w(\mathrm{S}) \cap \mathrm{T} = \{0\}\) , hence

\[
\operatorname{codim} _ {\mathrm{F}} (\operatorname{Im} (w)) \geqslant \dim (\mathrm{T}) - \operatorname{codim} _ {\mathrm{E}} (\mathrm{S}) = r + n + 1,
\]

this entails \(\operatorname{ind}(w) \geqslant r + 1\) and contradicts the hypothesis that v is adherent to B \(_{r}\) .

If \(r\) is an element of \(\mathbf{Z}\) such that \(B_r\) is nonempty, we have \(B_r = B\) since \(B_r\) is open and closed in \(B\) and \(B\) is connected. If \(B_r\) is empty for every \(r \in \mathbf{Z}\) , we have \(\text{ind}(v) = +\infty\) for all \(v \in B\) . The map \(v \mapsto \text{ind}(v)\) is therefore constant on \(B\) . Finally, for every \(v \in B\) , we have

\[
\begin{array}{c} \dim \operatorname{Coker} (v) = \operatorname{ind} (v) + \dim \operatorname{Ker} (v) \\ \leqslant \operatorname{ind} (u) + \dim \operatorname{Ker} (u) = \dim \operatorname{Coker} (u). \end{array}
\]

This concludes the proof.

COROLLARY 1. --- The functions defined by \(u \mapsto \dim \operatorname{Ker}(u)\) and by \(u \mapsto \dim \operatorname{Coker}(u)\) on \(\mathcal{Q}\mathcal{M}(\mathrm{E};\mathrm{F})\) are upper semicontinuous. The function \(u \mapsto \operatorname{ind}(u)\) is locally constant on \(\mathcal{Q}\mathcal{M}(\mathrm{E};\mathrm{F})\) .

COROLLARY 2. --- The function \(u \mapsto \dim \operatorname{Coker}(u)\) is locally constant on the set \(\mathcal{M}(\mathrm{E};\mathrm{F})\) of strict injective morphisms from E into F.

Indeed, one has \(\dim \operatorname{Coker}(u) = \operatorname{ind}(u)\) for \(u \in \mathcal{M}(\mathrm{E};\mathrm{F})\) .

Lemma 4. --- Let E and F be Banach spaces. For an element u of \(\mathcal{L}(\mathrm{E};\mathrm{F})\) to be a strict morphism from E into F, it is necessary and sufficient that \(^{t}u\) be a strict morphism from F' into E'. In this case, one has \(\dim \operatorname{Coker}(^{t}u)=\dim\operatorname{Ker}(u)\) and \(\dim\operatorname{Ker}(^{t}u)=\dim\operatorname{Coker}(u)\) .

For u to be a strict morphism, it is necessary and sufficient that \(^{t}u\) be one (EVT, IV, p. 29, cor. 3); in this case, the image of u is closed (EVT, IV, p. 28, th. 1) and the vector space \(\mathrm{Ker}(^{t}u)\) (resp. \(\mathrm{Coker}(^{t}u)\) ) is canonically isomorphic to the dual of the normed space \(\mathrm{Coker}(u)\) (resp. \(\mathrm{Ker}(u)\) ) according to EVT, IV, p. 27, prop. 2. The lemma follows.

Let E and F be Banach spaces. We denote by \(\mathcal{E}(\mathrm{E};\mathrm{F})\) the set of continuous surjective linear maps from E into F and by \(\mathcal{Q}\mathcal{E}(\mathrm{E};\mathrm{F})\) the set of continuous linear maps from E into F whose image has finite codimension. Every element of \(\mathcal{Q}\mathcal{E}(\mathrm{E};\mathrm{F})\) is a strict morphism with closed image (III, p. 52, lemma 6). It follows from lemma 4 that \(\mathcal{E}(\mathrm{E};\mathrm{F})\) and \(\mathcal{Q}\mathcal{E}(\mathrm{E};\mathrm{F})\) are respectively the inverse images of \(\mathcal{M}(\mathrm{F}^{\prime};\mathrm{E}^{\prime})\) and \(\mathcal{Q}\mathcal{M}(\mathrm{F}^{\prime};\mathrm{E}^{\prime})\) under the continuous map \(u\mapsto{}^{t}u\) from \(\mathcal{L}(\mathrm{E};\mathrm{F})\) into \(\mathcal{L}(\mathrm{F}^{\prime};\mathrm{E}^{\prime})\) .

PROPOSITION 13. --- Let E and F be Banach spaces. The sets \(\mathcal{E}(\mathrm{E};\mathrm{F})\) and \(\mathcal{Q}\mathcal{E}(\mathrm{E};\mathrm{F})\) are open in \(\mathcal{L}(\mathrm{E};\mathrm{F})\) . More precisely, if \(u\) is an element of \(\mathcal{Q}\mathcal{E}(\mathrm{E};\mathrm{F})\) and \(v\) is an element of \(\mathcal{L}(\mathrm{E};\mathrm{F})\) such that \(\| v - u\| < ((u))\) , we have \(v \in \mathcal{Q}\mathcal{E}(\mathrm{E};\mathrm{F})\) and

\(\dim \operatorname{Ker}(v) \leqslant \dim \operatorname{Ker}(u), \quad \dim \operatorname{Coker}(v) \leqslant \dim \operatorname{Coker}(u),\)

\[
\operatorname{ind} (v) = \operatorname{ind} (u).
\]

We saw above that \(^{t}u \in \mathcal{Q}\mathcal{M}(F';E')\) . Moreover, for every element \(v\) of \(\mathcal{L}(E;F)\) , we have \(\|^{t}v - ^{t}u\| = \|v - u\|\) and \(((^{t}u)) = ((u))\) (EVT, IV, p. 7, prop. 8 and III, p. 63, prop. 8). By prop. 12, it follows from these relations that if \(\|v - u\| < ((u))\) , then \(^{t}v\) belongs to \(\mathcal{Q}\mathcal{M}(F';E')\) and one has the inequalities

\[
\dim \operatorname{Ker} (^ {t} v) \leqslant \dim \operatorname{Ker} (^ {t} u), \quad \dim \operatorname{Coker} (^ {t} v) \leqslant \dim \operatorname{Coker} (^ {t} u)
\]

as well as the equality \(\mathrm{ind}(^{t}v)=\mathrm{ind}(^{t}u)\) . The proposition then follows from lemma 4.

COROLLARY 1. --- The functions defined by \(u \mapsto \dim \operatorname{Ker}(u)\) and by \(u \mapsto \dim \operatorname{Coker}(u)\) on \(\mathcal{Q}\mathcal{E}(\mathrm{E};\mathrm{F})\) are upper semicontinuous. The function \(u \mapsto \operatorname{ind}(u)\) is locally constant on \(\mathcal{Q}\mathcal{E}(\mathrm{E};\mathrm{F})\) .

COROLLARY 2. --- The function defined by \(u \mapsto \dim \operatorname{Ker}(u)\) is locally constant on \(\mathcal{E}(\mathrm{E};\mathrm{F})\) .

Indeed, one has \(\dim \operatorname{Ker}(u) = -\operatorname{ind}(u)\) for all \(u \in \mathcal{E}(\mathrm{E};\mathrm{F})\) .

